Triangulation from one vertex gives angle sum $180(n-2)^\circ$. Thus $n-2\le\dfrac{2026}{180}$, and the largest integer $n$ is $\boxed{13}$.
\begin{asy}
import graph;
usepackage("amsmath");
import geometry;
size(7cm);
pair[] p={(3.0,0.0),(2.65637,1.39417),(1.70419,2.46895),(0.36161,2.97813),(-1.06381,2.80505),(-2.24553,1.98937),(-2.91283,0.71795),(-2.91283,-0.71795),(-2.24553,-1.98937),(-1.06381,-2.80505),(0.36161,-2.97813),(1.70419,-2.46895),(2.65637,-1.39417)}; pair labelCentre=(0,0);for(pair labelVertex:p)labelCentre+=labelVertex;labelCentre/=p.length; draw(p[0]--p[1]--p[2]--p[3]--p[4]--p[5]--p[6]--p[7]--p[8]--p[9]--p[10]--p[11]--p[12]--cycle); dot(p[0]);label("$A$",p[0],unit(p[0]-labelCentre));dot(p[1]);label("$B$",p[1],unit(p[1]-labelCentre));dot(p[2]);label("$C$",p[2],unit(p[2]-labelCentre));dot(p[3]);label("$D$",p[3],unit(p[3]-labelCentre));dot(p[4]);label("$E$",p[4],unit(p[4]-labelCentre));dot(p[5]);label("$F$",p[5],unit(p[5]-labelCentre));dot(p[6]);label("$G$",p[6],unit(p[6]-labelCentre));dot(p[7]);label("$H$",p[7],unit(p[7]-labelCentre));dot(p[8]);label("$I$",p[8],unit(p[8]-labelCentre));dot(p[9]);label("$J$",p[9],unit(p[9]-labelCentre));dot(p[10]);label("$K$",p[10],unit(p[10]-labelCentre));dot(p[11]);label("$L$",p[11],unit(p[11]-labelCentre));dot(p[12]);label("$M$",p[12],unit(p[12]-labelCentre));
\end{asy}

